% =====================================================================
% paper/manuscript.tex
%
% Phase 2 electron kinetics — short paper note
%
% This file is intentionally a one-page redirection note rather than a
% full manuscript. The primary Phase 2 deliverable is the technical
% report at report/technical_report.pdf. A separate exploratory
% identifiability side-study exists at
% supplementary/identifiability_M7/ with its own manuscript.
%
% Author: Zachariah Ngan, Illinois Plasma Institute
% =====================================================================
\documentclass[11pt,a4paper]{article}

\usepackage[margin=1.1in]{geometry}
\usepackage[utf8]{inputenc}
\usepackage[T1]{fontenc}
\usepackage{hyperref}
\usepackage{parskip}

\hypersetup{colorlinks=true, linkcolor=blue!55!black, urlcolor=blue!55!black}

\title{\textbf{Phase 2 Electron Kinetics Pipeline}\\[4pt]
       \large Paper note}
\author{Zachariah Ngan\\[2pt]
        \normalsize\itshape Illinois Plasma Institute\\
        \small\texttt{zngan@illinois.edu}}
\date{\today}

\begin{document}
\maketitle

\section*{About this file}

This \texttt{paper/} directory is \emph{not} the primary deliverable of
the Phase 2 electron-kinetics work for the Illinois Plasma Institute
SF$_6$ ICP modelling project. It is retained in the repository for
completeness, but a reader interested in Phase 2 itself should go
directly to the technical report.

\subsection*{Where the actual Phase 2 deliverable lives}

\begin{itemize}
  \item \textbf{Primary technical report:} \texttt{report/technical\_report.pdf}.
        Organised as Tier 1 (BOLSIG$^+$ lookup and Maxwellian
        comparison), Tier 2 (supervised MLP surrogate and the
        \texttt{get\_rates\_pinn()} production API), Tier 3 (Monte
        Carlo collision module and three workplan-specified
        validation cases). Matches the April 2026 workplan
        structure one-to-one.
  \item \textbf{Production code:} \texttt{tier1\_bolsig/},
        \texttt{tier2\_pinn/}, \texttt{tier3\_picmcc/} at the
        repository root. Each directory has its own README and
        runnable scripts.
  \item \textbf{Slide deck for supervisor presentation:}
        \texttt{slides/slides.pdf}.
  \item \textbf{Top-level orientation:} \texttt{README.md} at the
        repository root explains which directory contains what and
        how to reproduce the results.
\end{itemize}

\subsection*{Where the exploratory identifiability side-study lives}

An earlier analytical side-study on two-anchor etch-rate calibration
identifiability exists at
\texttt{supplementary/identifiability\_M7/}. That work is a separate
analytical note that addresses a different question from Phase 2 (it
asks what two-anchor calibration constrains; Phase 2 asks how much the
EEDF replacement moves the fluorine profile). It was built in an
earlier session before the supervisor's workplan was loaded into
context.

The side-study's main result --- that two etch-rate anchors constrain
only the product $\mathcal{I} = \beta\,k_{\mathrm{diss}}\,n_e(P_H)$
--- is internally consistent and potentially publishable on its own as
a short methodological note. It is \emph{not} the Phase 2 deliverable
and should not be cited as such. Its own manuscript
(\texttt{manuscript\_or\_note.tex}), figures, and reframing README are
preserved in \texttt{supplementary/identifiability\_M7/}.

\subsection*{Why this file exists as a stub}

When the Phase 2 repository was finalised against the supervisor's
workplan, the main question was whether the earlier exploratory
paper should be deleted, kept as the primary narrative, or moved to
supplementary material. The honest answer is none of the three:

\begin{itemize}
  \item Deleting it would discard internally-consistent analytical
        work that may still have value as a standalone methodological
        note.
  \item Keeping it as the primary narrative would misrepresent the
        Phase 2 deliverable, which is a three-tier production
        kinetics pipeline, not an identifiability theorem.
  \item Moving it silently to supplementary without leaving a pointer
        at \texttt{paper/} would disorient a reader who landed there
        from an external reference.
\end{itemize}

The compromise is this short stub: the \texttt{paper/} directory still
exists and still compiles, but its contents are a clearly-labelled
redirection note that tells the reader where to look for the actual
Phase 2 work and where to look for the analytical side-study. The
figures from the original paper are preserved in
\texttt{paper/figures/} so that the supplementary manuscript at
\texttt{supplementary/identifiability\_M7/} can share them without
duplication.

\vspace{2em}
\begin{center}
\textit{For the Phase 2 deliverable proper, open
\texttt{report/technical\_report.pdf}.}
\end{center}

\end{document}
